% Zone „Das kennst du schon“ – aus bank/lineare-funktionen/zone.jsonl (zusammenbau v0.1)
\einheitenkopf[zone][Kennst du schon]{Das kennst du schon}
\setcounter{aufgabe}{0}
%% TODO Grundvorstellungs-Aufgabe als erste Hauptnummer der Zone (2.8) fehlt in der Bank

%% TODO Titel: Ich-kann-Satz fehlt in der Bank (Platzhalter: Kettenname) – Nr. 1
%% TODO Anweisung: ein Satz über der ersten Teilaufgabe fehlt in der Bank – Nr. 1
\begin{aufgabe}{Koordinaten lesen und eintragen, 4 Quadranten, (x|y)-Reihenfolge}
\swz{Punkt A – welche Koordinaten? Lies ab.}{\begin{ksys}[xmin=-5,xmax=5,ymin=-5,ymax=5,ablesen] \punkt{4}{3}{A} \end{ksys}}
\swz{Punkt C – welche Koordinaten? Lies ab.}{\begin{ksys}[xmin=-5,xmax=5,ymin=-5,ymax=5,ablesen] \punkt{-3}{-4}{C} \end{ksys}}
\end{aufgabe}

%% TODO Titel: Ich-kann-Satz fehlt in der Bank (Platzhalter: Kettenname) – Nr. 2
%% TODO Anweisung: ein Satz über der ersten Teilaufgabe fehlt in der Bank – Nr. 2
\begin{aufgabe}{Proportionale Zuordnung erkennen und hochrechnen (Dreisatz)}
%% TODO Darstellung für schwach fehlt in der Bank (lineare-funktionen-zone-f2-v1; Abschnitt „Für schwache Schüler“)
\swz{Preis für 5 kg? 2 kg Äpfel kosten 4 €. Wie viel kosten 5 kg?}{}
%% TODO Darstellung für schwach fehlt in der Bank (lineare-funktionen-zone-f2-v3; Abschnitt „Für schwache Schüler“)
\swz{Preis für 3 Liter? 4 Liter Saft kosten 6 €. Wie viel kosten 3 Liter?}{}
\end{aufgabe}

%% TODO Titel: Ich-kann-Satz fehlt in der Bank (Platzhalter: Kettenname) – Nr. 3
%% TODO Anweisung: ein Satz über der ersten Teilaufgabe fehlt in der Bank – Nr. 3
\begin{aufgabe}{Negative Zahlen multiplizieren und dividieren}
%% TODO Darstellung für schwach fehlt in der Bank (lineare-funktionen-zone-f3-v1; Abschnitt „Für schwache Schüler“)
\swz{Rechne: $(-3) \cdot 4$}{}
%% TODO Darstellung für schwach fehlt in der Bank (lineare-funktionen-zone-f3-v3; Abschnitt „Für schwache Schüler“)
\swz{Rechne: $(-1{,}5) \cdot (-4)$}{}
\end{aufgabe}

%% TODO Titel: Ich-kann-Satz fehlt in der Bank (Platzhalter: Kettenname) – Nr. 4
%% TODO Anweisung: ein Satz über der ersten Teilaufgabe fehlt in der Bank – Nr. 4
\begin{aufgabe}{Brüche als Steigung (½, −¾), Bruch mal ganze Zahl}
%% TODO Darstellung für schwach fehlt in der Bank (lineare-funktionen-zone-f4-v1; Abschnitt „Für schwache Schüler“)
\swz{Rechne: $\frac{1}{2} \cdot 8$}{}
%% TODO Darstellung für schwach fehlt in der Bank (lineare-funktionen-zone-f4-v3; Abschnitt „Für schwache Schüler“)
\swz{Rechne: $-\frac{3}{4} \cdot 8$}{}
\end{aufgabe}

%% TODO Titel: Ich-kann-Satz fehlt in der Bank (Platzhalter: Kettenname) – Nr. 5
%% TODO Anweisung: ein Satz über der ersten Teilaufgabe fehlt in der Bank – Nr. 5
\begin{aufgabe}{Lineare Gleichung zweischrittig lösen (null gleich einem zweischrittigen Term)}
%% TODO Darstellung für schwach fehlt in der Bank (lineare-funktionen-zone-f5-v1; Abschnitt „Für schwache Schüler“)
\swz{$\text{Löse: $0 = 2x - 8$}$}{}
%% TODO Darstellung für schwach fehlt in der Bank (lineare-funktionen-zone-f5-v3; Abschnitt „Für schwache Schüler“)
\swz{$\text{Löse: $0 = -4x + 12$}$}{}
\end{aufgabe}

%% TODO Titel: Ich-kann-Satz fehlt in der Bank (Platzhalter: Kettenname) – Nr. 6
%% TODO Anweisung: ein Satz über der ersten Teilaufgabe fehlt in der Bank – Nr. 6
\begin{aufgabe}{Terme zusammenfassen}
%% TODO Darstellung für schwach fehlt in der Bank (lineare-funktionen-zone-f6-v1; Abschnitt „Für schwache Schüler“)
\swz{Fasse zusammen: $3x + 4x$}{}
%% TODO Darstellung für schwach fehlt in der Bank (lineare-funktionen-zone-f6-v3; Abschnitt „Für schwache Schüler“)
\swz{Fasse zusammen: $2x + 5 - 4x + 1$}{}
\end{aufgabe}

%% TODO Titel: Ich-kann-Satz fehlt in der Bank (Platzhalter: Kettenname) – Nr. 7
%% TODO Anweisung: ein Satz über der ersten Teilaufgabe fehlt in der Bank – Nr. 7
\begin{aufgabe}{Negative Zahlen multiplizieren und dividieren – Fehler finden}
\swz[3]{Ich finde den Fehler: Mia rechnet $-3 \cdot (-5) + 2 = -13$. Benenne den Fehler und rechne richtig.}{}
\end{aufgabe}

%% TODO Titel: Ich-kann-Satz fehlt in der Bank (Platzhalter: Kettenname) – Nr. 8
%% TODO Anweisung: ein Satz über der ersten Teilaufgabe fehlt in der Bank – Nr. 8
\begin{aufgabe}{Negative Zahlen multiplizieren und dividieren – selbst rechnen}
%% TODO Darstellung für schwach fehlt in der Bank (lineare-funktionen-zone-f3-v6; Abschnitt „Für schwache Schüler“)
\swz{Rechne: $-6 \cdot (-4) + 5$}{}
\end{aufgabe}
